\documentclass[11pt]{article}
\usepackage[a4paper,margin=1in]{geometry}
\usepackage{amsmath,amssymb,amsthm,microtype}
\newtheorem{theorem}{Theorem}[section]
\newtheorem{corollary}[theorem]{Corollary}
\theoremstyle{remark}\newtheorem{remark}[theorem]{Remark}
\title{Endpoint Hardy Membership and a Logarithmic Mismatch Barrier}
\author{Bin Wang}\date{July 2026}
\begin{document}\maketitle
\begin{abstract}
At the deterministic coefficient radius
$\rho_*=1.4267874838\ldots$, the actual complement-to-anchor mismatch has a
sharp qualitative split: it belongs to $H^2(\rho_*)$ for every fixed noise,
but it never belongs to $H^\infty(\rho_*)$.  Endpoint $H^2$ convergence would
still close the weighted prefix at $R=1.4$, with conversion constant
$4.992111\ldots$.  However the exact odd anchors and the complement mass cap
force an endpoint norm lower bound of order $1/\sqrt{\log M_\sigma}$, hence
at least order $1/\sqrt{\log(1/\sigma)}$.  This barrier tends to zero, so it
does not prove endpoint convergence or nonconvergence.
\end{abstract}

\section{Endpoint functions}
Let
\[
 \rho_*=q_*^{-1}=0.85\lambda,\qquad q=\frac12,\qquad
 g_\sigma(z)=\sum_{n\ge2}
 \frac{\tau_{\sigma,n}-a_n}{n}z^n.
\]
For the modulus complement,
$|\tau_{\sigma,n}|\le M_\sigma q^{n-2}$ with finite
$M_\sigma\le\sigma^{-1}$.  The deterministic anchors satisfy
$|a_n|<48q_*^n$ and, for odd $n$,
\[
 a_n=\frac{q_*^n}{1+\lambda^{-n}}.
\]

\section{Membership theorem}
\begin{theorem}\label{thm:membership}
For every fixed $\sigma>0$,
\[
 g_\sigma\in H^2(\rho_*),\qquad
 g_\sigma\notin H^\infty(\rho_*).
\]
\end{theorem}
\begin{proof}
The noisy logarithm is bounded at $\rho_*$ because
$q\rho_*=q/q_*<1$ and its coefficients have the geometric mass envelope.
The target logarithm belongs to $H^2(\rho_*)$ because
\[
 \sum_{n\ge2}\frac{|a_n|^2\rho_*^{2n}}{n^2}
 \le48^2\sum_{n\ge2}\frac1{n^2}<\infty.
\]
Thus their difference belongs to $H^2$.

The odd part of the target logarithm has positive radial values
\[
 \sum_{\substack{n\ge3\\n\ {\rm odd}}}
 \frac{r^n}{n(1+\lambda^{-n})},\qquad 0<r<1,
\]
which diverge as $r\uparrow1$.  Hence the target is not bounded on the
endpoint disk.  Subtracting the bounded noisy logarithm cannot make it
bounded, proving the $H^\infty$ exclusion.
\end{proof}

\section{Endpoint conversion and lower barrier}
Put $R=1.4$ and $x=R/\rho_*$.  Every endpoint $H^2$ mismatch satisfies
\[
 \sum_{n\ge2}\frac{|\tau_{\sigma,n}-a_n|R^n}{n}
 \le\|g_\sigma\|_{H^2(\rho_*)}
 \frac{x^2}{\sqrt{1-x^2}},
\]
where
\[
 \frac{x^2}{\sqrt{1-x^2}}
 =4.992111068649647\ldots.
\]

\begin{theorem}\label{thm:barrier}
Let $M>0$ and let $N(M)$ be the least odd $N\ge3$ satisfying
\[
 Mq^{N-2}\le\frac14q_*^N.
\]
Then every modulus-capped complement of squared mass $M$ obeys
\[
 \boxed{
 \|g\|_{H^2(\rho_*)}
 \ge\frac1{\sqrt{32N(M)}}.}
\]
Consequently, for constants $C,M_0>0$,
\[
 \|g\|_{H^2(\rho_*)}
 \ge\frac{C}{\sqrt{\log(e+M)}}
 \qquad(M\ge M_0).
\]
For the actual complement mass $M_\sigma\le\sigma^{-1}$,
\[
 \|g_\sigma\|_{H^2(\rho_*)}
 \ge\frac{C'}{\sqrt{\log(e+1/\sigma)}}.
\]
\end{theorem}
\begin{proof}
For every odd $n\ge N(M)$, the cap remains below $q_*^n/4$.  Since
$a_n\ge q_*^n/2$,
\[
 |\tau_n-a_n|\ge\frac14q_*^n.
\]
After scaling by $\rho_*^n/n$, every such odd coefficient contributes at
least $1/(4n)$ in modulus.  Therefore
\[
 \|g\|_{H^2(\rho_*)}^2
 \ge\frac1{16}
 \sum_{j\ge0}\frac1{(N+2j)^2}
 \ge\frac1{32N}.
\]
The defining inequality gives
$N(M)=\log M/\log(q_*/q)+O(1)$.  Since
$M\mapsto1/\sqrt{\log(e+M)}$ is decreasing, the archived upper cap
$M_\sigma\le\sigma^{-1}$ gives the final statement; bounded masses only
strengthen it.
\end{proof}

\section{Protocol and exact boundary}
The computation reports the exact forced odd cutoff and certified lower
bound for $M=10^4,10^8,10^{16}$, together with a separately named normalized
logarithmic scale.  Infinite model $\sum_{n>N}n^{-2}$ tails are represented
by rigorous integral bounds rather than a finite truncation.

\begin{remark}
The lower bound tends to zero.  It rules out endpoint $H^2$ decay faster
than the displayed logarithmic scale but proves neither endpoint convergence
nor divergence.
The endpoint $H^\infty$ route is unavailable because the actual mismatch is
not in that space.  Gates A--E remain false/open.
\end{remark}
\end{document}
